\begin{tikzpicture}
    \begin{axis}[
        width=0.95\linewidth,
        height=0.6\linewidth,
        xlabel={Zeit / \si{\second}},
        ylabel={Amplitude / \si{\volt}},
        % --- Grid-Optimierung ---
        grid=both,
        major grid style={dashed, gray!20}, % Hell und gestrichelt
        minor tick num=0,                   % Entfernt die feinen Zwischenlinien
        %xtick distance={0.1},             % Optional: Erzwingt Linien alle 0.1s
            % ------------------------
        legend cell align=left,
        legend pos=south east,
        tick label style={font=\footnotesize},
        xmin=0, xmax=0.43,
            %empty cells=keep,
    ]
        % Hintergrund (Oszillogramm)
        \addplot[color={rgb,255:red,10;green,85;blue,140}, opacity=0.15, forget plot, line width=0.3pt]
            table[x=time, y=V_hp, col sep=comma]{Bilder/sprung_osz_hptp.csv};
            
        \addplot[color={rgb,255:red,195;green,5;blue,35}, opacity=0.15, forget plot, line width=0.3pt]
            table[x=time, y=V_tp, col sep=comma]{Bilder/sprung_osz_hptp.csv};

            % Vordergrund (Hüllkurven)
        \addplot[color={rgb,255:red,10;green,85;blue,140}, thick, smooth]
            table[x=time, y=amp_hp, col sep=comma]{Bilder/sprung_env_hptp.csv};
        \addlegendentry{Hüllkurve HP}

        \addplot[color={rgb,255:red,195;green,5;blue,35}, thick, smooth]
            table[x=time, y=amp_tp, col sep=comma]{Bilder/sprung_env_hptp.csv};
        \addlegendentry{Hüllkurve TP}

            % Markierungen
            %\draw[black, dashed, thick] (axis cs:0.01,\pgfkeysvalueof{/pgfplots/ymin}) -- (axis cs:0.01,\pgfkeysvalueof{/pgfplots/ymax});
            %\draw[black, dashed, thick] (axis cs:0.16,\pgfkeysvalueof{/pgfplots/ymin}) -- (axis cs:0.16,\pgfkeysvalueof{/pgfplots/ymax});
            
            %\addlegendimage{black, dashed, thick}
            %\addlegendentry{Frequenzsprung}
        
    \end{axis} 
\end{tikzpicture}